\documentclass[11pt,a4paper]{article}

% Packages
\usepackage[utf8]{inputenc}
\usepackage[T1]{fontenc}
\usepackage{amsmath,amssymb,amsfonts}
\usepackage{mathtools}
\usepackage{bm}
\usepackage{geometry}
\usepackage{hyperref}
\usepackage{cleveref}
\usepackage{booktabs}
\usepackage{array}
\usepackage{enumitem}
\usepackage{tikz}
\usepackage{pgfplots}
\usepackage{fancyvrb}
\usepackage{xcolor}
\usepackage{tcolorbox}
\usepackage{caption}

\geometry{margin=1in}
\pgfplotsset{compat=1.18}

% Custom colors
\definecolor{implnote}{RGB}{40,80,40}
\definecolor{implbg}{RGB}{235,245,235}
\definecolor{dimcolor}{RGB}{0,60,120}
\definecolor{evobg}{RGB}{240,240,255}

% Custom environments
\newtcolorbox{implbox}[1][]{
  colback=implbg, colframe=implnote,
  fonttitle=\bfseries, title={Implementation Note},
  #1
}
\newtcolorbox{evobox}[1][]{
  colback=evobg, colframe=blue!40!black,
  fonttitle=\bfseries, title={Architectural Evolution},
  #1
}
\newtcolorbox{dimbox}[1][]{
  colback=blue!3, colframe=dimcolor,
  fonttitle=\bfseries, title={Concrete Dimensions},
  #1
}

% Notation shortcuts
\newcommand{\R}{\mathbb{R}}
\newcommand{\softmax}{\operatorname{softmax}}
\newcommand{\LayerNorm}{\operatorname{LayerNorm}}
\newcommand{\RMSNorm}{\operatorname{RMSNorm}}
\newcommand{\ReLU}{\operatorname{ReLU}}
\newcommand{\GELU}{\operatorname{GELU}}
\newcommand{\Swish}{\operatorname{Swish}}
\newcommand{\SwiGLU}{\operatorname{SwiGLU}}
\newcommand{\FFN}{\operatorname{FFN}}
\newcommand{\Attention}{\operatorname{Attention}}
\newcommand{\MultiHead}{\operatorname{MultiHead}}
\newcommand{\Concat}{\operatorname{Concat}}
\newcommand{\dmodel}{d_{\text{model}}}
\newcommand{\dff}{d_{\text{ff}}}
\newcommand{\dk}{d_k}
\newcommand{\dv}{d_v}

\title{\textbf{The Architecture of Large Language Models}\\[0.5em]
\large From the Original Transformer to Modern Decoder-Only Models}
\author{LLM Theory Reference}
\date{April 2026}

\begin{document}
\maketitle

\begin{abstract}
This document provides a top-down architectural reference for the Transformer and its
evolution into modern large language models. Starting from the original encoder-decoder
Transformer \cite{vaswani2017}, we trace the development through GPT-1 \cite{radford2018}
to contemporary decoder-only architectures like LLaMA \cite{touvron2023}. Each component
is presented with full mathematical formulations, matrix dimensions, and implementation
context. The goal is understanding \emph{what} each component does, \emph{why} it exists,
and \emph{how} modern models have refined it.
\end{abstract}

\tableofcontents
\newpage

%=============================================================================
\section{The Original Transformer}
\label{sec:transformer}
%=============================================================================

The Transformer \cite{vaswani2017} was introduced in 2017 as a sequence-to-sequence model
for machine translation. Its key innovation was replacing recurrence entirely with
\emph{attention mechanisms}, allowing the model to relate any two positions in a sequence
in constant time.

\subsection{Why Attention Replaced Recurrence}

Prior architectures (RNNs, LSTMs) processed sequences one step at a time, creating a
fundamental bottleneck: information from early tokens had to survive through every
intermediate step to reach later tokens. The Transformer eliminates this by letting every
position attend directly to every other position.

The computational trade-offs per layer are \cite{vaswani2017}:

\begin{center}
\begin{tabular}{lcc}
\toprule
\textbf{Layer Type} & \textbf{Complexity per Layer} & \textbf{Sequential Operations} \\
\midrule
Self-Attention & $O(n^2 \cdot d)$ & $O(1)$ \\
Recurrent & $O(n \cdot d^2)$ & $O(n)$ \\
\bottomrule
\end{tabular}
\end{center}

\noindent where $n$ is the sequence length and $d$ is the representation dimension.
Self-attention is more expensive per layer when $n > d$, but its $O(1)$ sequential
operations (fully parallelizable) make it dramatically faster on modern hardware. For
typical LLM configurations ($d = 4096$, $n \leq 8192$), both terms are comparable, and
the parallelism advantage dominates.

\subsection{High-Level Architecture}

The original Transformer is an \emph{encoder-decoder} model:

\begin{itemize}[nosep]
  \item \textbf{Encoder} (left): Processes the input sequence. Consists of $N$ identical
    layers, each containing a self-attention sub-layer followed by a feed-forward network.
    The encoder sees the entire input at once (bidirectional).
  \item \textbf{Decoder} (right): Generates the output sequence one token at a time.
    Also $N$ layers, but each layer has three sub-layers: \emph{masked} self-attention
    (can only attend to earlier positions), encoder-decoder cross-attention (attends to
    encoder output), and a feed-forward network.
\end{itemize}

Both stacks use residual connections and layer normalization around each sub-layer.
The base model uses $N = 6$ layers in each stack, with $\dmodel = 512$
\cite{vaswani2017}.

\begin{dimbox}
\textbf{Original Transformer (base):}
$N = 6$ layers, $\dmodel = 512$, $h = 8$ heads,
$\dk = \dv = \dmodel / h = 64$, $\dff = 2048$.
Total parameters: 65M.
\end{dimbox}


%=============================================================================
\section{Tokenization}
\label{sec:tokenization}
%=============================================================================

Before a model can process text, it must convert characters into discrete integer tokens.
The dominant method is \textbf{Byte Pair Encoding} (BPE) \cite{sennrich2016}, adapted from
a data compression algorithm.

\subsection{The BPE Algorithm}

BPE builds a vocabulary of subword units by iteratively merging the most frequent
adjacent pairs:

\begin{tcolorbox}[colback=white, colframe=black!50, title={Algorithm: BPE vocabulary construction \cite{sennrich2016}}]
\begin{enumerate}[nosep]
  \item Initialize vocabulary $\mathcal{V}$ with all individual characters (or bytes)
  \item Represent each word in the training corpus as a sequence of characters
  \item \textbf{For} $i = 1$ to $k$ (desired number of merges):
  \begin{enumerate}[nosep,label=\alph*.]
    \item Count all adjacent symbol pairs across the corpus
    \item Find the most frequent pair $(a, b)$
    \item Merge every occurrence of $(a, b)$ into a new symbol $ab$
    \item Add $ab$ to vocabulary $\mathcal{V}$
  \end{enumerate}
  \item \textbf{Return} $\mathcal{V}$
\end{enumerate}
\end{tcolorbox}

The only hyperparameter is the number of merge operations $k$, which directly controls
the final vocabulary size: $|\mathcal{V}| = |\text{base characters}| + k$.

\textbf{Example.} Starting from the corpus \{``low'', ``lowest'', ``newer'', ``wider''\},
BPE might learn merges: \texttt{l o} $\to$ \texttt{lo}, then \texttt{lo w} $\to$
\texttt{low}, then \texttt{e r} $\to$ \texttt{er}. After these merges, the unseen word
``lower'' segments as \texttt{low er} --- two known tokens \cite{sennrich2016}.

\subsection{Vocabulary Sizes in Practice}

\begin{center}
\begin{tabular}{lrl}
\toprule
\textbf{Model} & \textbf{Vocab Size} & \textbf{Tokenizer} \\
\midrule
Original Transformer \cite{vaswani2017} & 37,000 & BPE (shared source/target) \\
GPT-1 \cite{radford2018} & 40,000 & BPE \\
GPT-3 \cite{brown2020} & 50,257 & BPE \\
LLaMA \cite{touvron2023} & 32,000 & SentencePiece BPE (byte-fallback) \\
\bottomrule
\end{tabular}
\end{center}

\begin{implbox}
LLaMA uses SentencePiece \cite{touvron2023}, which operates on raw byte sequences
rather than Unicode characters. This provides a \emph{byte-fallback} mechanism: any byte
sequence can be tokenized, even if it contains characters unseen during training. All
numbers are split into individual digits. This means the vocabulary is truly closed ---
there are no ``unknown token'' failures at inference time.
\end{implbox}


%=============================================================================
\section{Embeddings}
\label{sec:embeddings}
%=============================================================================

Each token index $t \in \{0, 1, \ldots, V-1\}$ is mapped to a dense vector via the
\textbf{token embedding matrix}:
\begin{equation}
  W_e \in \R^{V \times \dmodel}
\end{equation}
where $V$ is the vocabulary size and $\dmodel$ is the model dimension. Looking up token
$t$ simply selects the $t$-th row of $W_e$.

For a sequence of tokens $(t_1, t_2, \ldots, t_n)$, the embedding step produces a matrix:
\begin{equation}
  X_{\text{embed}} = \begin{pmatrix} W_e[t_1] \\ W_e[t_2] \\ \vdots \\ W_e[t_n] \end{pmatrix}
  \in \R^{n \times \dmodel}
\end{equation}

\subsection{Weight Tying}
\label{sec:weight-tying}

A key architectural detail: in both the original Transformer \cite{vaswani2017} and GPT-1
\cite{radford2018}, the \emph{same} weight matrix $W_e$ is used for three purposes:

\begin{enumerate}[nosep]
  \item Source (encoder) input embeddings
  \item Target (decoder) input embeddings
  \item The output projection that converts hidden states back to vocabulary logits
    (see \Cref{sec:output-head})
\end{enumerate}

In decoder-only models (which have no encoder), points 1 and 2 collapse, so $W_e$ is
shared between the input embedding and the output projection:
\begin{equation}
  \text{logits} = h_n \cdot W_e^\top
\end{equation}
where $h_n \in \R^{\dmodel}$ is the final hidden state. This halves the parameters
devoted to vocabulary (which can be substantial --- for LLaMA-7B with $V = 32000$ and
$\dmodel = 4096$, the embedding matrix alone is $32000 \times 4096 \approx 131$M
parameters).

\subsection{Embedding Scaling}

The original Transformer scales embeddings by $\sqrt{\dmodel}$ before adding positional
encodings \cite{vaswani2017}:
\begin{equation}
  X = \sqrt{\dmodel} \cdot X_{\text{embed}} + X_{\text{pos}}
\end{equation}
where $X_{\text{pos}} \in \R^{n \times \dmodel}$ is the positional encoding matrix
(\Cref{sec:positional-encoding}).
This compensates for the fact that embedding vectors tend to have small magnitudes
(roughly unit variance per dimension), while positional encodings have magnitudes on the
order of 1. The scaling ensures the token identity signal is not overwhelmed by positional
information. GPT-1 and LLaMA do not use this scaling factor.


%=============================================================================
\section{Positional Encoding}
\label{sec:positional-encoding}
%=============================================================================

Self-attention is \emph{permutation-invariant}: if you shuffle the input tokens,
the attention computation produces the same outputs (just shuffled). The model has no
inherent sense of token order. Positional encodings inject sequence-position information
so the model can distinguish ``the cat sat on the mat'' from ``mat the on sat cat the.''

Three major approaches have been used, representing an evolution in how position is
encoded.

\subsection{Sinusoidal Positional Encoding (2017)}

The original Transformer \cite{vaswani2017} adds fixed sinusoidal functions to the
embeddings. For position $\text{pos}$ and dimension $i$:
\begin{align}
  PE_{(\text{pos}, 2i)}   &= \sin\!\left(\frac{\text{pos}}{10000^{2i/\dmodel}}\right) \\[4pt]
  PE_{(\text{pos}, 2i+1)} &= \cos\!\left(\frac{\text{pos}}{10000^{2i/\dmodel}}\right)
\end{align}
Each dimension oscillates at a different frequency, from $2\pi$ (dimension 0) to
$10000 \cdot 2\pi$ (last dimension). The intuition: position is encoded as a binary-like
code where each dimension captures a different ``bit'' of the position at a different
time scale.

A key property: for any fixed offset $k$, $PE_{\text{pos}+k}$ can be expressed as a
linear function of $PE_{\text{pos}}$. This means the model can learn to attend to
relative positions using linear projections.

These encodings are \emph{added} to the token embeddings (not concatenated), so they share
the same $\dmodel$-dimensional space.

\subsection{Learned Positional Embeddings (2018)}

GPT-1 \cite{radford2018} replaced the fixed sinusoidal functions with a \emph{learned}
position embedding matrix:
\begin{equation}
  W_p \in \R^{n_{\max} \times \dmodel}
\end{equation}
where $n_{\max}$ is the maximum sequence length (512 for GPT-1). The position embedding
for position $\text{pos}$ is simply the $\text{pos}$-th row of $W_p$, trained jointly with
the rest of the model:
\begin{equation}
  h_0 = X_{\text{embed}} \cdot W_e + W_p
\end{equation}

This gives the model full freedom to learn whatever positional representation is most
useful, but fixes the maximum context length at training time.

\subsection{Rotary Position Embedding --- RoPE (2021)}
\label{sec:rope}

RoPE \cite{su2021} takes a fundamentally different approach: instead of \emph{adding}
position to the embeddings, it \emph{rotates} the query and key vectors by
position-dependent angles. This makes the dot product between query at position $m$ and
key at position $n$ depend only on their relative distance $(m - n)$.

For a vector $\bm{x} \in \R^d$, RoPE groups dimensions into pairs and applies a 2D
rotation to each pair. The rotation angle for pair $i$ at position $m$ is:
\begin{equation}
  \theta_i = m \cdot \omega_i, \quad \text{where } \omega_i = 10000^{-2i/d}
\end{equation}
where $m$ is the token position in the sequence, $i \in \{0, 1, \ldots, d/2 - 1\}$ is
the dimension pair index, and $d$ is the head dimension ($\dk$).

The rotation is applied as:
\begin{equation}
  f(\bm{x}, m) =
  \begin{pmatrix}
    x_0 \cos(m\omega_0) - x_1 \sin(m\omega_0) \\
    x_0 \sin(m\omega_0) + x_1 \cos(m\omega_0) \\
    x_2 \cos(m\omega_1) - x_3 \sin(m\omega_1) \\
    x_2 \sin(m\omega_1) + x_3 \cos(m\omega_1) \\
    \vdots
  \end{pmatrix}
\end{equation}
where $x_0, x_1, \ldots$ are the scalar components of $\bm{x}$.

The key mathematical property is that for rotated queries $\bm{q}_m = f(\bm{q}, m)$ and
keys $\bm{k}_n = f(\bm{k}, n)$:
\begin{equation}
  \bm{q}_m^\top \bm{k}_n = \bm{q}^\top R_{m-n} \bm{k}
\end{equation}
where $R_{m-n}$ is a rotation matrix depending only on the relative position $(m - n)$.
The dot product --- and thus the attention weight --- naturally encodes relative position
without any explicit relative position bias.

\begin{evobox}
\textbf{RoPE} is applied only to the query $Q$ and key $K$ vectors, \emph{not} to values
$V$ \cite{su2021}. It is applied at every layer, not just the first. This is the
positional encoding used by LLaMA \cite{touvron2023} and most modern LLMs.
\end{evobox}


%=============================================================================
\section{Attention Mechanism}
\label{sec:attention}
%=============================================================================

Attention is the core operation of the Transformer. Conceptually, it allows each position
in a sequence to ``look at'' every other position and aggregate information based on
relevance.

\subsection{Scaled Dot-Product Attention}

Given queries $Q$, keys $K$, and values $V$, attention computes \cite{vaswani2017}:
\begin{equation}
\label{eq:attention}
  \Attention(Q, K, V) = \softmax\!\left(\frac{QK^\top}{\sqrt{\dk}}\right) V
\end{equation}
where $Q, K \in \R^{n \times \dk}$ are the query and key matrices, $V \in \R^{n \times \dv}$
is the value matrix, $n$ is the sequence length, and $\dk$ is the key/query dimension
(the scaling factor).

Step by step, for a single query vector $\bm{q}$ attending to keys $K$ and values $V$:

\begin{enumerate}[nosep]
  \item \textbf{Compute similarity scores:} $\bm{s} = \bm{q} K^\top \in \R^n$. Each
    score $s_j = \bm{q} \cdot \bm{k}_j$ measures how relevant position $j$ is.
  \item \textbf{Scale:} $\bm{s} \leftarrow \bm{s} / \sqrt{\dk}$. The scaling prevents
    the dot products from growing large in magnitude with $\dk$, which would push the
    softmax into regions of extremely small gradients. Without scaling, for random unit-variance
    vectors, $\text{Var}(\bm{q} \cdot \bm{k}) = \dk$, so dividing by $\sqrt{\dk}$
    restores unit variance \cite{vaswani2017}.
  \item \textbf{Softmax:} $\bm{a} = \softmax(\bm{s}) \in \R^n$. Converts scores to a
    probability distribution (attention weights). Each $a_j \geq 0$ and $\sum_j a_j = 1$.
  \item \textbf{Weighted sum:} $\text{output} = \sum_j a_j \bm{v}_j$. The output is a
    weighted combination of value vectors, where the weights are the attention scores.
\end{enumerate}

In matrix form for all $n$ positions simultaneously:
\begin{equation}
  Q, K \in \R^{n \times \dk}, \quad
  V \in \R^{n \times \dv}, \quad
  QK^\top \in \R^{n \times n}, \quad
  \Attention(Q, K, V) \in \R^{n \times \dv}
\end{equation}

The $n \times n$ attention matrix is the computational bottleneck --- it's why
self-attention is $O(n^2)$ in sequence length.

\subsection{Multi-Head Attention}

Rather than performing a single attention function with $\dmodel$-dimensional keys,
values, and queries, multi-head attention runs $h$ attention operations in parallel, each
on a different learned linear projection \cite{vaswani2017}:
\begin{align}
  \MultiHead(Q, K, V) &= \Concat(\text{head}_1, \ldots, \text{head}_h) \, W^O \\
  \text{where } \text{head}_i &= \Attention(Q W_i^Q, \; K W_i^K, \; V W_i^V)
\end{align}

The projection matrices are:
\begin{equation}
  W_i^Q \in \R^{\dmodel \times \dk}, \quad
  W_i^K \in \R^{\dmodel \times \dk}, \quad
  W_i^V \in \R^{\dmodel \times \dv}, \quad
  W^O \in \R^{h\dv \times \dmodel}
\end{equation}

With $\dk = \dv = \dmodel / h$, the total computation is similar to single-head attention
with full dimensionality, but the model can attend to information from different
representation subspaces at different positions simultaneously.

\begin{dimbox}
\textbf{Original Transformer:} $\dmodel = 512$, $h = 8$, so $\dk = \dv = 64$. Each head
operates on 64-dimensional projections.\\[4pt]
\textbf{LLaMA-7B:} $\dmodel = 4096$, $h = 32$, so $\dk = \dv = 128$. With RoPE applied
to $Q$ and $K$ at every layer.
\end{dimbox}

\subsection{Three Types of Attention}
\label{sec:attention-types}

The original encoder-decoder Transformer uses three distinct attention patterns
\cite{vaswani2017}:

\begin{enumerate}[nosep]
  \item \textbf{Encoder self-attention:} $Q$, $K$, $V$ all come from the encoder.
    Every position can attend to every other position (bidirectional). The attention
    matrix is dense.

  \item \textbf{Decoder masked self-attention:} $Q$, $K$, $V$ all come from the decoder.
    Position $i$ can only attend to positions $\leq i$ (autoregressive / causal). This
    is enforced by setting the upper-triangular entries of $QK^\top$ to $-\infty$ before
    the softmax, so they become zero:
    \begin{equation}
      \text{mask}_{ij} =
      \begin{cases}
        0 & \text{if } j \leq i \\
        -\infty & \text{if } j > i
      \end{cases}
    \end{equation}
    where $i$ is the query position and $j$ is the key position.

  \item \textbf{Encoder-decoder cross-attention:} $Q$ comes from the decoder, $K$ and $V$
    come from the encoder output. This is how the decoder ``reads'' the input. Each
    decoder position can attend to all encoder positions.
\end{enumerate}

In decoder-only models (GPT, LLaMA), only type 2 exists. Cross-attention and bidirectional
self-attention are removed entirely (\Cref{sec:decoder-only}).


%=============================================================================
\section{Feed-Forward Network}
\label{sec:ffn}
%=============================================================================

Each Transformer layer contains a \textbf{position-wise feed-forward network} (FFN) ---
a two-layer fully-connected network applied independently and identically to each position
in the sequence \cite{vaswani2017}. While attention mixes information \emph{across}
positions, the FFN transforms the representation \emph{at} each position.

\subsection{Original FFN (ReLU)}

\begin{equation}
  \FFN(x) = \ReLU(x W_1 + b_1) W_2 + b_2
  \label{eq:ffn-relu}
\end{equation}
where $x \in \R^{\dmodel}$ is the hidden state at a single position,
$W_1 \in \R^{\dmodel \times \dff}$, $W_2 \in \R^{\dff \times \dmodel}$ are weight matrices,
$b_1 \in \R^{\dff}$, $b_2 \in \R^{\dmodel}$ are bias vectors, and
$\ReLU(z) = \max(0, z)$.

The inner dimension $\dff$ is typically $4 \times \dmodel$. In the original Transformer,
$\dmodel = 512$ and $\dff = 2048$, so the FFN ``expands'' the representation to 4x width,
applies a nonlinearity, then ``compresses'' back. This expansion is thought to provide
capacity for the model to process and transform features.

\subsection{GELU Activation (GPT-1)}

GPT-1 \cite{radford2018} replaced ReLU with the \textbf{Gaussian Error Linear Unit}
(GELU):
\begin{equation}
  \GELU(x) = x \cdot \Phi(x)
\end{equation}
where $\Phi(x)$ is the cumulative distribution function of the standard normal
distribution. Unlike ReLU, which hard-clips at zero, GELU applies a smooth,
input-dependent gate. Values near zero are partially attenuated rather than fully zeroed.

\begin{equation}
  \FFN_{\text{GELU}}(x) = \GELU(x W_1) W_2
  \label{eq:ffn-gelu}
\end{equation}
where $x$, $W_1$, and $W_2$ are as defined in \Cref{eq:ffn-relu} (without bias terms).

(GPT-1 also drops the bias terms, following the T5 convention \cite{shazeer2020}.)

\subsection{SwiGLU (LLaMA and Modern Models)}

Shazeer \cite{shazeer2020} introduced Gated Linear Unit (GLU) variants into the
Transformer FFN. The key idea is to use a \emph{gating mechanism}: one linear projection
produces the content, another produces a gate that controls what passes through.

\textbf{SwiGLU} combines the Swish activation function with a gated linear unit:
\begin{equation}
  \FFN_{\SwiGLU}(x, W, V, W_2) = \left(\Swish_1(xW) \otimes xV \right) W_2
  \label{eq:ffn-swiglu}
\end{equation}
where $x \in \R^{\dmodel}$ is the input hidden state,
$W, V \in \R^{\dmodel \times \dff}$ are the gate and up-projection weight matrices,
$W_2 \in \R^{\dff \times \dmodel}$ is the down-projection weight matrix,
$\otimes$ is element-wise multiplication, $\Swish_\beta(x) = x \cdot \sigma(\beta x)$,
and $\sigma$ is the sigmoid function. When $\beta = 1$, this is also called SiLU.

GLU variants use \emph{three} weight matrices ($W$, $V$, $W_2$) instead of the
standard two ($W_1$, $W_2$). To keep the parameter count and computation equivalent, the
inner dimension $\dff$ is reduced by a factor of $\frac{2}{3}$ \cite{shazeer2020}:
\begin{equation}
  \dff^{\text{SwiGLU}} = \frac{2}{3} \cdot 4\dmodel
\end{equation}

\begin{dimbox}
\textbf{LLaMA-7B:} $\dmodel = 4096$, $\dff = \frac{2}{3} \times 4 \times 4096 = 10922
\approx 11008$ (rounded to a multiple of 256 for hardware efficiency) \cite{touvron2023}.
Compare to the standard $4 \times 4096 = 16384$ that a non-gated FFN would use.
\end{dimbox}

\begin{evobox}
\textbf{FFN activation evolution:}
\begin{itemize}[nosep]
  \item 2017: ReLU \cite{vaswani2017} --- hard zero below 0
  \item 2018: GELU \cite{radford2018} --- smooth, probabilistic gating
  \item 2020: SwiGLU \cite{shazeer2020} --- gated architecture with Swish activation
\end{itemize}
SwiGLU consistently achieves better perplexity than ReLU and GELU in matched-parameter
comparisons \cite{shazeer2020}, and is used in LLaMA \cite{touvron2023} and PaLM
\cite{chowdhery2022}. The
original paper offers ``no explanation as to why these architectures seem to work;
we attribute their success, as all else, to divine benevolence'' \cite{shazeer2020}.
\end{evobox}


%=============================================================================
\section{Normalization and Residual Connections}
\label{sec:norm-residual}
%=============================================================================

Every sub-layer (attention, FFN) in the Transformer is wrapped in a residual connection
and a normalization step. The placement of normalization relative to the sub-layer has
significant implications for training dynamics.

\subsection{Residual Connections}

Each sub-layer output is added to its input via a skip connection:
\begin{equation}
  \text{output} = x + \text{SubLayer}(x)
\end{equation}
where $x$ is the sub-layer input and $\text{SubLayer}(\cdot)$ is either the attention
or FFN function.

Residual connections serve two critical purposes:
\begin{enumerate}[nosep]
  \item \textbf{Gradient flow:} They create a direct path for gradients to flow from
    the loss back to early layers, mitigating the vanishing gradient problem in deep
    networks.
  \item \textbf{Incremental refinement:} Each layer's job becomes learning a
    \emph{correction} to its input rather than computing the entire representation from
    scratch. This makes optimization easier.
\end{enumerate}

\subsection{Layer Normalization}

Layer normalization \cite{xiong2020} normalizes across the feature dimension for each
individual position and sample:
\begin{equation}
  \LayerNorm(\bm{v}) = \gamma \cdot \frac{\bm{v} - \mu}{\sigma} + \beta
\end{equation}
where $\bm{v} \in \R^d$ is the input vector ($d$ is the feature dimension, i.e.\ $\dmodel$),
$\mu = \frac{1}{d}\sum_{k=1}^{d} v_k$ is the mean,
$\sigma = \sqrt{\frac{1}{d}\sum_{k=1}^{d}(v_k - \mu)^2}$ is the standard deviation,
and $\gamma, \beta \in \R^d$ are learned scale and bias parameters.

\subsection{Post-LN vs.\ Pre-LN}
\label{sec:post-pre-ln}

The placement of layer normalization has major implications for training stability
\cite{xiong2020}:

\textbf{Post-LN} (original Transformer \cite{vaswani2017}):
\begin{align}
  x' &= \LayerNorm(x + \text{SubLayer}(x))
\end{align}
where $x$ is the layer input and $x'$ is the layer output.
The sub-layer output is added to the input first, \emph{then} normalized.

\textbf{Pre-LN} (GPT and modern models):
\begin{align}
  x' &= x + \text{SubLayer}(\LayerNorm(x))
\end{align}
The input is normalized first, \emph{then} passed through the sub-layer, then added back.
Pre-LN also requires a \emph{final} layer normalization before the output projection.

Xiong et al.\ \cite{xiong2020} proved via mean field theory that these placements have
different gradient behaviors at initialization:

\begin{center}
\begin{tabular}{lcc}
\toprule
\textbf{Variant} & \textbf{Output-layer gradient scale} & \textbf{Warmup needed?} \\
\midrule
Post-LN & $O(d\sqrt{\ln d})$, independent of $L$ & Yes (critical) \\
Pre-LN  & $O\!\left(d\sqrt{\frac{\ln d}{L}}\right)$, decreases with depth & No \\
\bottomrule
\end{tabular}
\end{center}

\noindent where $d$ is the model dimension ($\dmodel$) and $L$ is the number of layers.

In Post-LN, gradients near the output layer are large regardless of model depth,
making training unstable without careful learning rate warmup. In Pre-LN, gradients
scale inversely with $\sqrt{L}$ (number of layers), keeping them well-behaved even for
deep models. This is why Pre-LN enables training without warmup and is used by virtually
all modern LLMs.

\subsection{RMSNorm}
\label{sec:rmsnorm}

RMSNorm \cite{zhang2019} simplifies layer normalization by removing the mean-centering
step:
\begin{equation}
  \RMSNorm(\bm{v}) = \gamma \cdot \frac{\bm{v}}{\text{RMS}(\bm{v})},
  \quad \text{where } \text{RMS}(\bm{v}) = \sqrt{\frac{1}{d}\sum_{k=1}^{d} v_k^2}
\end{equation}
where $\bm{v} \in \R^d$ is the input vector, $d$ is the feature dimension,
$v_k$ are the scalar components of $\bm{v}$, and $\gamma \in \R^d$ is a learned
scale parameter.

By dropping the mean subtraction and the bias term $\beta$, RMSNorm provides:
\begin{itemize}[nosep]
  \item Comparable model quality to standard LayerNorm
  \item 7--64\% reduction in computation time \cite{zhang2019}
\end{itemize}

LLaMA \cite{touvron2023} uses Pre-LN with RMSNorm --- this is now the standard
configuration for modern LLMs.

\begin{evobox}
\textbf{Normalization evolution:}
\begin{itemize}[nosep]
  \item 2017: Post-LN with standard LayerNorm \cite{vaswani2017} --- requires warmup
  \item 2018: Pre-LN with standard LayerNorm \cite{radford2018} --- no warmup needed
  \item 2023: Pre-LN with RMSNorm \cite{touvron2023} --- no warmup, faster computation
\end{itemize}
\end{evobox}


%=============================================================================
\section{The Decoder-Only Shift}
\label{sec:decoder-only}
%=============================================================================

Modern LLMs (GPT \cite{radford2018, brown2020}, LLaMA \cite{touvron2023}, Mistral
\cite{jiang2023}, Claude, etc.) are all \emph{decoder-only} Transformers. This section
traces why and how the field converged on this architecture.

\subsection{From Encoder-Decoder to Decoder-Only}

GPT-1 \cite{radford2018} made the critical architectural choice: take the Transformer
decoder, remove the encoder-decoder cross-attention sub-layer, and use the resulting model
for language modeling. The architecture reduces to:

\begin{align}
  h_0 &= X_{\text{embed}} \cdot W_e + W_p \\
  h_l &= \text{transformer\_block}(h_{l-1}) \quad \forall\, l \in [1, L] \\
  P(t) &= \softmax(h_L \cdot W_e^\top)
\end{align}
where $h_0 \in \R^{n \times \dmodel}$ is the initial hidden state,
$X_{\text{embed}}$ is the token index matrix, $W_e$ is the embedding matrix,
$W_p \in \R^{n_{\max} \times \dmodel}$ is the learned position embedding,
$h_l$ is the hidden state after layer $l$, $L$ is the total number of layers,
and $P(t) \in \R^V$ is the output probability distribution over the vocabulary.

Each transformer block consists of just two sub-layers:
\begin{enumerate}[nosep]
  \item Masked multi-head self-attention
  \item Position-wise feed-forward network
\end{enumerate}

\subsection{Causal Masking}

The defining feature of decoder-only models is the \textbf{causal mask}: token at position
$i$ can only attend to tokens at positions $\leq i$. This is implemented by adding a mask
to the attention scores before the softmax:

\begin{equation}
  \Attention(Q, K, V) = \softmax\!\left(\frac{QK^\top}{\sqrt{\dk}} + M \right) V
\end{equation}

where $M \in \R^{n \times n}$ has $M_{ij} = 0$ if $j \leq i$ and $M_{ij} = -\infty$
if $j > i$.

This enforces autoregressive generation: the prediction for position $i$ depends only on
tokens $1, \ldots, i$, matching the language modeling objective
$P(t_i \mid t_1, \ldots, t_{i-1})$.

\subsection{Why Decoder-Only Won}

The encoder-decoder structure was designed for sequence-to-sequence tasks (translation),
where you have a complete input and generate a complete output. For general language
modeling and generation:

\begin{itemize}[nosep]
  \item There is no separate ``input'' to encode --- the model simply continues a sequence.
  \item The decoder alone is sufficient for the autoregressive objective.
  \item Removing the encoder and cross-attention cuts the architecture roughly in half,
    simplifying both the model and the training pipeline.
  \item The same model serves both understanding and generation --- no separate encoder
    needed.
\end{itemize}

\subsection{Architectural Timeline}

\begin{center}
\begin{tabular}{llp{7.5cm}}
\toprule
\textbf{Model} & \textbf{Year} & \textbf{Key changes from Transformer} \\
\midrule
Transformer \cite{vaswani2017} & 2017 &
  Encoder-decoder, Post-LN, sinusoidal PE, ReLU FFN \\
GPT-1 \cite{radford2018} & 2018 &
  Decoder-only, Pre-LN\footnotemark, learned PE, GELU, BPE 40K \\
GPT-3 \cite{brown2020} & 2020 &
  Scaled to 175B params (96 layers, $d{=}12288$, 96 heads). Alternating dense and
  locally-banded sparse attention in some layers \\
LLaMA \cite{touvron2023} & 2023 &
  Pre-RMSNorm, SwiGLU ($\dff = \frac{2}{3} \cdot 4d$), RoPE at every layer, no bias
  terms, SentencePiece BPE. Sizes 7B--65B, trained on public data only \\
\bottomrule
\end{tabular}
\end{center}
\footnotetext{GPT-1's Figure 1 shows Layer Norm placed before each sub-layer (the Pre-LN
arrangement), which differs from the original Transformer's Post-LN.}


%=============================================================================
\section{Output Head}
\label{sec:output-head}
%=============================================================================

The output head converts the final hidden state into a probability distribution over the
vocabulary.

\subsection{Final Layer Normalization}

In Pre-LN models, the residual stream after the last layer has not been normalized (since
Pre-LN normalizes \emph{before} each sub-layer). A final layer normalization is therefore
applied before the output projection \cite{xiong2020}:
\begin{equation}
  h_{\text{final}} = \RMSNorm(h_L) \quad \text{(or LayerNorm in older models)}
\end{equation}
where $h_L \in \R^{n \times \dmodel}$ is the output of the last Transformer layer ($L$ layers total).

\subsection{Projection to Logits}

The normalized hidden state is projected to vocabulary-sized logits:
\begin{equation}
  \bm{z} = h_{\text{final}} \cdot W_e^\top \in \R^V
\end{equation}
where $W_e \in \R^{V \times \dmodel}$ is the (shared) embedding matrix
(\Cref{sec:weight-tying}). Each element $z_j$ is the unnormalized log-probability
(logit) for vocabulary token $j$.

\subsection{Softmax}

The logits are converted to probabilities:
\begin{equation}
  P(t = j \mid t_1, \ldots, t_{i-1}) = \frac{e^{z_j}}{\sum_{k=0}^{V-1} e^{z_k}}
\end{equation}
where $z_j$ is the logit for token $j$, and $V$ is the vocabulary size.

During training, the target is the next token $t_i$, and the loss is the cross-entropy:
\begin{equation}
  \mathcal{L} = -\log P(t_i \mid t_1, \ldots, t_{i-1})
\end{equation}
summed over all positions in the sequence.

During inference, a token is selected from this distribution (via greedy decoding, top-$k$
sampling, nucleus sampling, etc. --- these are inference-time strategies outside the scope
of this document).


%=============================================================================
\section{Putting It Together}
\label{sec:full-forward-pass}
%=============================================================================

Let us trace a complete forward pass through a modern decoder-only Transformer, using
LLaMA-7B dimensions \cite{touvron2023} as a concrete example.

\begin{dimbox}
\textbf{LLaMA-7B hyperparameters:}
$\dmodel = 4096$, $L = 32$ layers, $h = 32$ heads, $\dk = \dv = 128$,
$\dff = 11008$, $V = 32{,}000$. Total: $\sim$6.7B parameters.
\end{dimbox}

\subsection{Step-by-Step Forward Pass}

Given an input sequence of $n$ token indices $(t_1, \ldots, t_n)$:

\medskip

\noindent\textbf{Step 1: Token Embedding}
\begin{equation}
  X = W_e[t_1, \ldots, t_n] \in \R^{n \times 4096}
\end{equation}
Look up each token index in $W_e \in \R^{32000 \times 4096}$.

\medskip

\noindent\textbf{Step 2: Repeat for each of $L = 32$ layers.} Let $h^{(0)} = X$.

For layer $l = 1, \ldots, 32$:

\medskip

\noindent\textbf{Step 2a: Pre-RMSNorm (attention)}
\begin{equation}
  \hat{h} = \RMSNorm(h^{(l-1)}) \in \R^{n \times 4096}
\end{equation}

\noindent\textbf{Step 2b: Multi-head self-attention with RoPE}
\begin{align}
  Q &= \hat{h} \, W^Q \in \R^{n \times 4096} \quad \text{(reshaped to $n \times 32 \times 128$)} \\
  K &= \hat{h} \, W^K \in \R^{n \times 4096} \quad \text{(reshaped to $n \times 32 \times 128$)} \\
  V &= \hat{h} \, W^V \in \R^{n \times 4096} \quad \text{(reshaped to $n \times 32 \times 128$)}
\end{align}
Apply RoPE rotations to $Q$ and $K$ based on position.

For each head $i$, compute:
\begin{equation}
  \text{head}_i = \softmax\!\left(\frac{Q_i K_i^\top}{\sqrt{128}} + M\right) V_i
  \in \R^{n \times 128}
\end{equation}
where $M$ is the causal mask. Concatenate all 32 heads and project:
\begin{equation}
  A = \Concat(\text{head}_1, \ldots, \text{head}_{32}) \, W^O \in \R^{n \times 4096}
\end{equation}

\noindent\textbf{Step 2c: Residual connection (attention)}
\begin{equation}
  h' = h^{(l-1)} + A
\end{equation}

\noindent\textbf{Step 2d: Pre-RMSNorm (FFN)}
\begin{equation}
  \hat{h}' = \RMSNorm(h') \in \R^{n \times 4096}
\end{equation}

\noindent\textbf{Step 2e: SwiGLU FFN}
\begin{equation}
  F = \left(\Swish(\hat{h}' W_{\text{gate}}) \otimes \hat{h}' W_{\text{up}}\right) W_{\text{down}}
  \in \R^{n \times 4096}
\end{equation}
where $W_{\text{gate}}, W_{\text{up}} \in \R^{4096 \times 11008}$ and
$W_{\text{down}} \in \R^{11008 \times 4096}$.

\noindent\textbf{Step 2f: Residual connection (FFN)}
\begin{equation}
  h^{(l)} = h' + F
\end{equation}

\medskip

\noindent\textbf{Step 3: Final RMSNorm}
\begin{equation}
  h_{\text{final}} = \RMSNorm(h^{(32)}) \in \R^{n \times 4096}
\end{equation}

\noindent\textbf{Step 4: Output projection}
\begin{equation}
  \text{logits} = h_{\text{final}} \cdot W_e^\top \in \R^{n \times 32000}
\end{equation}

\noindent\textbf{Step 5: Softmax (per position)}
\begin{equation}
  P(t_{i+1} \mid t_1, \ldots, t_i) = \softmax(\text{logits}[i]) \in \R^{32000}
\end{equation}

\subsection{Parameter Count Breakdown}

\begin{center}
\begin{tabular}{lrrr}
\toprule
\textbf{Component} & \textbf{Params per layer} & \textbf{$\times$ 32 layers} & \textbf{Total} \\
\midrule
Attention ($W^Q, W^K, W^V, W^O$) & $4 \times 4096^2$ & & 2.1B \\
FFN ($W_{\text{gate}}, W_{\text{up}}, W_{\text{down}}$) & $3 \times 4096 \times 11008$ & & 4.3B \\
RMSNorm ($\gamma$ $\times$ 2) & $2 \times 4096$ & & 8.4M \\
\midrule
All layers & & & 6.4B \\
Token embedding $W_e$ & & & 131M \\
Final RMSNorm & & & 4K \\
\midrule
\textbf{Total} & & & $\sim$\textbf{6.7B} \\
\bottomrule
\end{tabular}
\end{center}

Note: the FFN accounts for roughly two-thirds of per-layer parameters, and attention
accounts for one-third. This ratio holds across model sizes.


%=============================================================================
% References
%=============================================================================
\newpage
\begin{thebibliography}{9}

\bibitem{vaswani2017}
  A.~Vaswani, N.~Shazeer, N.~Parmar, J.~Uszkoreit, L.~Jones, A.~N.~Gomez,
  \L.~Kaiser, and I.~Polosukhin,
  ``Attention Is All You Need,''
  \emph{Advances in Neural Information Processing Systems 30 (NeurIPS)}, 2017.
  \\\url{https://arxiv.org/abs/1706.03762}

\bibitem{radford2018}
  A.~Radford, K.~Narasimhan, T.~Salimans, and I.~Sutskever,
  ``Improving Language Understanding by Generative Pre-Training,''
  OpenAI, 2018.
  \\\url{https://cdn.openai.com/research-covers/language-unsupervised/language_understanding_paper.pdf}

\bibitem{brown2020}
  T.~B.~Brown, B.~Mann, N.~Ryder, et al.,
  ``Language Models are Few-Shot Learners,''
  \emph{Advances in Neural Information Processing Systems 33 (NeurIPS)}, 2020.
  \\\url{https://arxiv.org/abs/2005.14165}

\bibitem{touvron2023}
  H.~Touvron, T.~Lavril, G.~Izacard, et al.,
  ``LLaMA: Open and Efficient Foundation Language Models,''
  arXiv preprint arXiv:2302.13971, 2023.
  \\\url{https://arxiv.org/abs/2302.13971}

\bibitem{sennrich2016}
  R.~Sennrich, B.~Haddow, and A.~Birch,
  ``Neural Machine Translation of Rare Words with Subword Units,''
  \emph{Proceedings of the 54th Annual Meeting of the ACL}, 2016.
  \\\url{https://arxiv.org/abs/1508.07909}

\bibitem{su2021}
  J.~Su, Y.~Lu, S.~Pan, B.~Wen, and Y.~Liu,
  ``RoFormer: Enhanced Transformer with Rotary Position Embedding,''
  arXiv preprint arXiv:2104.09864, 2021.
  \\\url{https://arxiv.org/abs/2104.09864}

\bibitem{shazeer2020}
  N.~Shazeer,
  ``GLU Variants Improve Transformer,''
  arXiv preprint arXiv:2002.05202, 2020.
  \\\url{https://arxiv.org/abs/2002.05202}

\bibitem{xiong2020}
  R.~Xiong, Y.~Yang, D.~He, et al.,
  ``On Layer Normalization in the Transformer Architecture,''
  \emph{Proceedings of the 37th International Conference on Machine Learning (ICML)}, 2020.
  \\\url{https://arxiv.org/abs/2002.04745}

\bibitem{zhang2019}
  B.~Zhang and R.~Sennrich,
  ``Root Mean Square Layer Normalization,''
  \emph{Advances in Neural Information Processing Systems 32 (NeurIPS)}, 2019.
  \\\url{https://arxiv.org/abs/1910.07467}

\end{thebibliography}

\end{document}
